\documentclass[11pt,a4paper]{article}
\usepackage{fontspec,amsmath,amssymb,amsthm,microtype}
\setmainfont{lmroman10-regular.otf}[BoldFont=lmroman10-bold.otf,ItalicFont=lmroman10-italic.otf,BoldItalicFont=lmroman10-bolditalic.otf]
\usepackage{xcolor}
\usepackage[most]{tcolorbox}
\usepackage{mathrsfs}
\providecommand{\lt}{<}
\providecommand{\gt}{>}
\setmonofont{DejaVuSansMono.ttf}[BoldFont=DejaVuSansMono-Bold.ttf,ItalicFont=DejaVuSansMono-Oblique.ttf,BoldItalicFont=DejaVuSansMono-BoldOblique.ttf,Scale=0.86]
\newfontfamily\leanfallback{STIXTwoMath-Regular.otf}[Scale=0.86]
\definecolor{leanrule}{gray}{0.62}
\definecolor{leanback}{gray}{0.962}
\definecolor{leantitle}{gray}{0.22}
\definecolor{leanmuted}{gray}{0.38}
\newtcolorbox{leanbox}[1]{enhanced,breakable,colback=leanback,colframe=leanrule,
  boxrule=0.4pt,arc=0.6mm,left=2.4mm,right=2.4mm,top=1.8mm,bottom=1.8mm,
  before skip=0.8em,after skip=1em,fonttitle=\small\bfseries,coltitle=leantitle,
  colbacktitle=leanback,titlerule=0.3pt,title={#1}}
\newlength{\leancharwidth}
\newenvironment{leancode}{\par\smallskip\begingroup\ttfamily\footnotesize\raggedright
  \settowidth{\leancharwidth}{M}\setlength{\parskip}{0pt}\setlength{\parindent}{0pt}
  \frenchspacing}{\par\endgroup\smallskip}
\newcommand{\leanline}[2]{\hangindent=\dimexpr#1\leancharwidth+4\leancharwidth\relax
  \hangafter=1\noindent\hspace*{#1\leancharwidth}#2\par}
\newcommand{\leansource}[1]{\par{\footnotesize\color{leanmuted}#1}\par}
\newcommand{\leanlabel}[1]{\par\smallskip\noindent{\small\bfseries #1}\par}
\usepackage[margin=25mm]{geometry}
\usepackage[unicode,hidelinks]{hyperref}
\hypersetup{pdfauthor={LangevinLean Project}}
\setlength{\emergencystretch}{3em}
\setlength{\parskip}{0.4em}
\setlength{\parindent}{0pt}
\allowdisplaybreaks[2]
\raggedbottom
\title{\texorpdfstring{The exact scale-invariant step and logarithmic iteration count}{The exact scale-invariant step and logarithmic iteration count}}
\author{LangevinLean Project}
\date{}
\begin{document}\maketitle
Draft mathematical exposition; these notes may contain mistakes.


This section proves scalar accuracy arithmetic. For \(\alpha,\beta,d,\varepsilon>0\), define

\[h=\frac1\beta\min\left(1,\frac{\varepsilon}{12\sqrt d}\right),
\qquad K=\frac\beta\alpha\max\left(1,\frac{12\sqrt d}{\varepsilon}\right).\]


The constants \(6\) and \(12\), the maximum, minimum and logarithm are the literal quantities of Corollary 1.1. A convergence estimate for actual sampling laws is not established by this scalar argument.


\section{\texorpdfstring{The chosen step and its reciprocal}{The chosen step and its reciprocal}}


\textbf{Claim.} Put \(m=\max(1,12\sqrt d/\varepsilon)\). Then \(h=1/(\beta m)>0\), \(h\le1/\beta\), and \(\alpha hK=1\).




\textbf{Proof.} If \(\varepsilon\le12\sqrt d\), the minimum in the step is its second term and the maximum is \(12\sqrt d/\varepsilon\). Otherwise they both select their first terms. In either case multiplication verifies \(h=1/(\beta m)\). Positivity follows from \(\beta>0\), \(m\ge1\), and the upper bound follows by replacing the minimum by \(1\). Finally substitute \(K=\beta m/\alpha\) and cancel the positive factors.

\begin{leanbox}{Lean: \texttt{paperStep}, \texttt{iterationPrefactor}, \texttt{paperStep\_\allowbreak{}eq\_\allowbreak{}inv} and 3 more}
\leansource{\texttt{paperStep} \quad WassersteinMixing/\allowbreak{}Complexity.lean, lines 19–20}
\begin{leancode}
\leanline{0}{def\ paperStep\ (β\ d\ ε\ :\ ℝ)\ :\ ℝ\ :=}
\leanline{2}{(1\ /\ β)\ *\ min\ 1\ (ε\ /\ (12\ *\ Real.sqrt\ d))}
\end{leancode}
\leansource{\texttt{iterationPrefactor} \quad WassersteinMixing/\allowbreak{}Complexity.lean, lines 23–24}
\begin{leancode}
\leanline{0}{def\ iterationPrefactor\ (α\ β\ d\ ε\ :\ ℝ)\ :\ ℝ\ :=}
\leanline{2}{(β\ /\ α)\ *\ max\ 1\ (12\ *\ Real.sqrt\ d\ /\ ε)}
\end{leancode}
\leansource{\texttt{paperStep\_\allowbreak{}eq\_\allowbreak{}inv} \quad WassersteinMixing/\allowbreak{}Complexity.lean, lines 27–40}
\begin{leancode}
\leanline{0}{theorem\ paperStep\_eq\_inv\ (β\ d\ ε\ :\ ℝ)\ (hβ\ :\ 0\ <\ β)\ (hd\ :\ 0\ <\ d)\ (hε\ :\ 0\ <\ ε)\ :}
\leanline{4}{paperStep\ β\ d\ ε\ =\ 1\ /\ (β\ *\ max\ 1\ (12\ *\ Real.sqrt\ d\ /\ ε))}
\end{leancode}
\leansource{\texttt{paperStep\_\allowbreak{}pos} \quad WassersteinMixing/\allowbreak{}Complexity.lean, lines 43–46}
\begin{leancode}
\leanline{0}{theorem\ paperStep\_pos\ (β\ d\ ε\ :\ ℝ)\ (hβ\ :\ 0\ <\ β)\ (hd\ :\ 0\ <\ d)\ (hε\ :\ 0\ <\ ε)\ :}
\leanline{4}{0\ <\ paperStep\ β\ d\ ε}
\end{leancode}
\leansource{\texttt{paperStep\_\allowbreak{}le\_\allowbreak{}inv\_\allowbreak{}beta} \quad WassersteinMixing/\allowbreak{}Complexity.lean, lines 49–52}
\begin{leancode}
\leanline{0}{theorem\ paperStep\_le\_inv\_beta\ (β\ d\ ε\ :\ ℝ)\ (hβ\ :\ 0\ <\ β)\ (\_hd\ :\ 0\ <\ d)\ (\_hε\ :\ 0\ <\ ε)\ :}
\leanline{4}{paperStep\ β\ d\ ε\ ≤\ 1\ /\ β}
\end{leancode}
\leansource{\texttt{step\_\allowbreak{}mul\_\allowbreak{}iterationPrefactor} \quad WassersteinMixing/\allowbreak{}Complexity.lean, lines 55–61}
\begin{leancode}
\leanline{0}{theorem\ step\_mul\_iterationPrefactor\ (α\ β\ d\ ε\ :\ ℝ)\ (hα\ :\ 0\ <\ α)}
\leanline{4}{(hβ\ :\ 0\ <\ β)\ (hd\ :\ 0\ <\ d)\ (hε\ :\ 0\ <\ ε)\ :}
\leanline{4}{(α\ *\ paperStep\ β\ d\ ε)\ *\ iterationPrefactor\ α\ β\ d\ ε\ =\ 1}
\end{leancode}
\leanlabel{How the Lean proof follows the argument}
\texttt{paper\allowbreak{}Step} and \texttt{iteration\allowbreak{}Prefactor} retain the exact minimum and maximum. \texttt{paper\allowbreak{}Step\_\allowbreak{}eq\_\allowbreak{}inv} checks the two branches explicitly. The positivity, stability and reciprocal-product lemmas clear only nonzero factors established from the positive hypotheses.
\end{leanbox}

\section{\texorpdfstring{Allocating half of the accuracy to the bias}{Allocating half of the accuracy to the bias}}


\textbf{Claim.} For \(\alpha>0\), \(\beta,d\ge0\),

\[\sqrt\alpha\sqrt{(\beta/\alpha)\beta d}=\beta\sqrt d.\]


Both sides are nonnegative and their squares agree after cancellation of \(\alpha\), which proves the identity. For the selected step and \(d>0\), it follows that

\[\sqrt\alpha\,6h\sqrt{(\beta/\alpha)\beta d}
=\frac{6\sqrt d}{m}\le\frac\varepsilon2.\]


The last inequality is precisely \(12\sqrt d\le\varepsilon m\), supplied by the second term in the maximum.

\begin{leanbox}{Lean: \texttt{sqrt\_\allowbreak{}curvature\_\allowbreak{}scale}, \texttt{paperStep\_\allowbreak{}bias\_\allowbreak{}le}}
\leansource{\texttt{sqrt\_\allowbreak{}curvature\_\allowbreak{}scale} \quad WassersteinMixing/\allowbreak{}Complexity.lean, lines 64–68}
\begin{leancode}
\leanline{0}{theorem\ sqrt\_curvature\_scale\ (α\ β\ d\ :\ ℝ)\ (hα\ :\ 0\ <\ α)\ (hβ\ :\ 0\ ≤\ β)\ (hd\ :\ 0\ ≤\ d)\ :}
\leanline{4}{Real.sqrt\ α\ *\ Real.sqrt\ ((β\ /\ α)\ *\ β\ *\ d)\ =\ β\ *\ Real.sqrt\ d}
\end{leancode}
\leansource{\texttt{paperStep\_\allowbreak{}bias\_\allowbreak{}le} \quad WassersteinMixing/\allowbreak{}Complexity.lean, lines 71–88}
\begin{leancode}
\leanline{0}{theorem\ paperStep\_bias\_le\ (α\ β\ d\ ε\ :\ ℝ)\ (hα\ :\ 0\ <\ α)\ (hβ\ :\ 0\ <\ β)}
\leanline{4}{(hd\ :\ 0\ <\ d)\ (hε\ :\ 0\ <\ ε)\ :}
\leanline{4}{Real.sqrt\ α\ *\ (6\ *\ paperStep\ β\ d\ ε\ *\ Real.sqrt\ ((β\ /\ α)\ *\ β\ *\ d))\ ≤\ ε\ /\ 2}
\end{leancode}
\leanlabel{How the Lean proof follows the argument}
\texttt{sqrt\_\allowbreak{}curvature\_\allowbreak{}scale} proves equality by equality of squares and nonnegativity. \texttt{paper\allowbreak{}Step\_\allowbreak{}bias\_\allowbreak{}le} substitutes this identity and the reciprocal step formula, then uses the maximum's lower bound to prove the half-budget inequality.
\end{leanbox}

\section{\texorpdfstring{The transient estimate and the complete scalar accuracy bound}{The transient estimate and the complete scalar accuracy bound}}


\textbf{Claim.} Let \(0<t\le1\), \(A\ge0\), \(\varepsilon>0\). If \(k\) is a natural number with \(\log(1+2A/\varepsilon)\le tk\), then

\[(1-t)^kA\le\varepsilon/2.\]


\textbf{Proof.} Since \(1-t\le e^{-t}\) and both are nonnegative, \((1-t)^k\le e^{-tk}\). The logarithmic hypothesis gives \(e^{-tk}\le(1+2A/\varepsilon)^{-1}\). Multiplication by \(A\ge0\) yields the result because \(A/(1+2A/\varepsilon)\le\varepsilon/2\). This includes \(A=0\) and \(t=1\).




Take \(t=\alpha h\), \(A=\sqrt\alpha R\), with \(R\ge0\) and \(0<\alpha\le\beta\). The chosen step gives \(t\le1\). If \(k\ge K\log(1+2\sqrt\alpha R/\varepsilon)\), then \(tK=1\) supplies the logarithmic hypothesis. Adding the transient and bias half-budgets proves

\[\sqrt\alpha\left((1-\alpha h)^kR+6h\sqrt{(\beta/\alpha)\beta d}\right)\le\varepsilon.\]\begin{leanbox}{Lean: \texttt{geometric\_\allowbreak{}transient\_\allowbreak{}le}, \texttt{paper\_\allowbreak{}complexity\_\allowbreak{}bound}}
\leansource{\texttt{geometric\_\allowbreak{}transient\_\allowbreak{}le} \quad WassersteinMixing/\allowbreak{}Complexity.lean, lines 92–114}
\begin{leancode}
\leanline{0}{theorem\ geometric\_transient\_le\ (t\ A\ ε\ :\ ℝ)\ (\_ht\ :\ 0\ <\ t)\ (ht1\ :\ t\ ≤\ 1)}
\leanline{4}{(hA\ :\ 0\ ≤\ A)\ (hε\ :\ 0\ <\ ε)\ (k\ :\ ℕ)}
\leanline{4}{(hk\ :\ Real.log\ (1\ +\ 2\ *\ A\ /\ ε)\ ≤\ t\ *\ k)\ :}
\leanline{4}{(1\ -\ t)\ \textasciicircum{}\ k\ *\ A\ ≤\ ε\ /\ 2}
\end{leancode}
\leansource{\texttt{paper\_\allowbreak{}complexity\_\allowbreak{}bound} \quad WassersteinMixing/\allowbreak{}Complexity.lean, lines 118–142}
\begin{leancode}
\leanline{0}{theorem\ paper\_complexity\_bound\ (α\ β\ d\ R\ ε\ :\ ℝ)\ (hα\ :\ 0\ <\ α)\ (hαβ\ :\ α\ ≤\ β)}
\leanline{4}{(hd\ :\ 0\ <\ d)\ (hR\ :\ 0\ ≤\ R)\ (hε\ :\ 0\ <\ ε)\ (k\ :\ ℕ)}
\leanline{4}{(hk\ :\ iterationPrefactor\ α\ β\ d\ ε\ *}
\leanline{6}{Real.log\ (1\ +\ 2\ *\ Real.sqrt\ α\ *\ R\ /\ ε)\ ≤\ (k\ :\ ℝ))\ :}
\leanline{4}{Real.sqrt\ α\ *\ ((1\ -\ α\ *\ paperStep\ β\ d\ ε)\ \textasciicircum{}\ k\ *\ R\ +}
\leanline{6}{6\ *\ paperStep\ β\ d\ ε\ *\ Real.sqrt\ ((β\ /\ α)\ *\ β\ *\ d))\ ≤\ ε}
\end{leancode}
\leanlabel{How the Lean proof follows the argument}
\texttt{geometric\_\allowbreak{}transient\_\allowbreak{}le} bounds the geometric factor by an exponential, uses monotonicity of the exponential to invert the logarithmic bound and checks the final quotient. \texttt{paper\_\allowbreak{}complexity\_\allowbreak{}bound} proves the admissible dimensionless step, invokes the reciprocal-product identity and adds the two half-budgets. No mixing hypothesis is hidden in these scalar statements.
\end{leanbox}

\end{document}
